% 单源最短路的三类结果。数据（经 Python 同步 Bellman--Ford 验算）：
% s->u 4, s->a 1, a->b 1, b->a -3, b->t 2, x->y 1, y->x -2, y->u 1。
% 结果：u=4（有限），a,b,t=-inf（可达负环 a->b->a，成本 -2），x,y=+inf（不可达）。
\begin{tikzpicture}[>=Stealth,
  v/.style={circle,draw=FigureStroke,fill=FigureBlueFill,line width=.8pt,minimum size=6.5mm,inner sep=0pt},
  vneg/.style={v,draw=FigureRed,fill=FigureRedFill},
  vinf/.style={v,draw=FigureMuted,fill=white,text=FigureMuted},
  arc/.style={->,draw=FigureStroke,line width=.8pt},
  w/.style={fill=white,inner sep=.8pt,font=\figurefont\fontsize{7.5}{9}\selectfont},
  note/.style={font=\FigureNoteFont\fontsize{7.5}{9}\selectfont}]
  \node[v] (s) at (0,0) {$s$};
  \node[v] (u) at (2,1.3) {$u$};
  \node[vneg] (a) at (2,-.8) {$a$};
  \node[vneg] (b) at (4,-.8) {$b$};
  \node[vneg] (t) at (6,-.8) {$t$};
  \node[vinf] (y) at (4,1.3) {$y$};
  \node[vinf] (x) at (6,1.3) {$x$};
  \draw[arc] (s)--node[w]{$4$}(u);
  \draw[arc] (s)--node[w]{$1$}(a);
  \draw[arc,draw=FigureRed] (a) to[bend left=28] node[w]{$1$}(b);
  \draw[arc,draw=FigureRed] (b) to[bend left=28] node[w]{$-3$}(a);
  \draw[arc] (b)--node[w]{$2$}(t);
  \draw[arc,draw=FigureMuted] (x) to[bend right=28] node[w]{$1$}(y);
  \draw[arc,draw=FigureMuted] (y) to[bend right=28] node[w]{$-2$}(x);
  \draw[arc,draw=FigureMuted] (y)--node[w]{$1$}(u);
  \node[note,text=FigureBlue,above=1mm of u] {$\delta(u)=4$};
  \node[note,text=FigureRed] at (3,-1.85) {负环$a\to b\to a$，成本$-2$};
  \node[note,text=FigureRed,above=1mm of t] {$\delta(t)=-\infty$};
  \node[note,text=FigureRed,left=1mm of a,yshift=-2.5mm] {$\delta(a)=-\infty$};
  \node[note,text=FigureRed,above=1mm of b] {$\delta(b)=-\infty$};
  \node[note,text=FigureMuted,above=1mm of y] {$\delta(y)=+\infty$};
  \node[note,text=FigureMuted,above=1mm of x] {$\delta(x)=+\infty$};
  \node[note,below=1mm of s] {$\delta(s)=0$};
\end{tikzpicture}
